% -------------------------------------------------------
%  Embedded Systems HW3 -- context.tex
%  آشنایی با Linux Embedded
% -------------------------------------------------------

%% ============================================================
\section{بخش تئوری}
%% ============================================================

\subsection{سؤال ۱: معماری کلی یک سیستم \lr{Linux Embedded}}

فرآیند راه‌اندازی یک سیستم \lr{Linux Embedded} از مجموعه‌ای از لایه‌ها تشکیل شده که هرکدام محیط اجرای لایه‌ی بعدی را فراهم می‌کنند:

\begin{itemize}
	\item \textbf{\lr{Bootloader}:} اولین کد اجراشده پس از روشن شدن برد (مانند \lr{U-Boot} یا \lr{Barebox}) است. \lr{Bootloader} معمولاً در چند مرحله اجرا می‌شود: مرحله‌ی اول (\lr{SPL}) که مستقیماً توسط سخت‌افزار از حافظه‌ی داخلی (\lr{ROM/SRAM}) خوانده می‌شود و فقط کار مقداردهی اولیه‌ی حداقلیِ \lr{DDR} و کنترلر ذخیره‌سازی را انجام می‌دهد، و مرحله‌ی دوم که وظایف اصلی را بر عهده می‌گیرد: مقداردهی اولیه‌ی کامل سخت‌افزار (کلاک‌ها، حافظه، پین‌مالتی‌پلکسینگ)، بارگذاری \lr{Linux Kernel} و \lr{Device Tree} از حافظه‌ی فلش/کارت حافظه به \lr{RAM}، تحویل آرگومان‌های خط فرمان کرنل (\lr{kernel command line})، و در نهایت پرش به نقطه‌ی ورود کرنل. به این ترتیب \lr{Bootloader} زمینه‌ی سخت‌افزاری و اطلاعات لازم برای اجرای کرنل را آماده می‌کند.

	\item \textbf{\lr{Linux Kernel}:} پس از دریافت کنترل از \lr{Bootloader}، کرنل ابتدا معماری پردازنده و \lr{Device Tree} را تفسیر کرده و بر اساس آن، درایورهای مربوط به سخت‌افزارهای توصیف‌شده (پردازنده، حافظه، گذرگاه‌ها، پیرامونی‌ها) را فعال می‌کند. کرنل زیرسیستم‌های حافظه (\lr{MMU}, \lr{paging})، زمان‌بند (\lr{scheduler})، شبکه و فایل‌سیستم مجازی را راه‌اندازی کرده و در پایان، \lr{Root File System} را \lr{mount} می‌کند و اولین پردازش کاربری (\lr{init}) را اجرا می‌کند. به این ترتیب کرنل، محیط اجرای نرم‌افزار کاربر را بر پایه‌ی سخت‌افزار فراهم می‌سازد.

	\item \textbf{\lr{Root File System}:} ساختار فایل‌ها و پوشه‌هایی است که کرنل پس از \lr{mount} کردن به آن دسترسی پیدا می‌کند و شامل تمام باینری‌ها، کتابخانه‌ها، فایل‌های پیکربندی و اسکریپت‌های راه‌اندازی سیستم است. بدون \lr{Root File System}، کرنل هیچ برنامه‌ی کاربری‌ای برای اجرا پیدا نمی‌کند و راه‌اندازی متوقف می‌شود (\lr{kernel panic}).

	\item \textbf{\lr{C Library}:} کتابخانه‌ای مانند \lr{glibc}، \lr{musl} یا \lr{uClibc} که رابط استاندارد بین برنامه‌های کاربر و \lr{system call}های کرنل را فراهم می‌کند. تقریباً تمام برنامه‌های کاربردی (از جمله خود \lr{init}) به این کتابخانه \lr{link} می‌شوند؛ بنابراین \lr{C Library} باید پیش از اجرای هر برنامه‌ی کاربری در \lr{Root File System} حاضر باشد.

	\item \textbf{\lr{init} / \lr{systemd}:} اولین پردازش کاربری (\lr{PID 1}) است که کرنل مستقیماً اجرا می‌کند. وظیفه‌ی آن راه‌اندازی و مدیریت چرخه‌ی حیات سایر سرویس‌ها و پردازش‌های سیستم است؛ به‌عنوان مثال \lr{systemd} با خواندن فایل‌های \lr{unit}، ترتیب و وابستگی سرویس‌ها را مشخص کرده و آن‌ها را (مانند شبکه، \lr{logging}، و برنامه‌های کاربر) راه‌اندازی می‌کند.

	\item \textbf{\lr{User Applications}:} برنامه‌های نهایی کاربر (سرویس‌ها، دیمن‌ها، برنامه‌های تعاملی) که توسط \lr{init}/\lr{systemd} یا به‌صورت دستی اجرا می‌شوند و منطق نهایی سیستم نهفته را پیاده‌سازی می‌کنند.
\end{itemize}

به‌طور خلاصه، هر لایه محیط اجرای لایه‌ی بعد را آماده می‌کند: \lr{Bootloader} سخت‌افزار را برای کرنل آماده می‌کند، کرنل \lr{Root File System} را برای \lr{init} در دسترس قرار می‌دهد، \lr{C Library} رابط اجرای برنامه‌ها را فراهم می‌سازد، و \lr{init}/\lr{systemd} در نهایت \lr{User Applications} را اجرا می‌کند.

%% ────────────────────────────────────────────────────────────
\subsection{سؤال ۲: محیط توسعه و \lr{Toolchain}}

\subsubsection*{الف) چرا \lr{Native Compilation} کافی نیست؟}

اگر برنامه را با \lr{g++} نصب‌شده روی لپ‌تاپ \lr{x86} کامپایل کنیم، کامپایلر به‌صورت پیش‌فرض کدی تولید می‌کند که:
\begin{itemize}
	\item برای مجموعه‌دستورالعمل (\lr{Instruction Set Architecture}) پردازنده‌ی \lr{x86} است، نه \lr{ARM}؛ پردازنده‌ی \lr{ARM} قادر به اجرای دستورالعمل‌های ماشین \lr{x86} نیست.
	\item در برابر \lr{C Library} و کتابخانه‌های سیستمی نصب‌شده روی \lr{Host} (لپ‌تاپ) \lr{link} می‌شود، که ممکن است روی \lr{Target} (برد \lr{ARM}) اصلاً وجود نداشته باشند یا نسخه‌ی متفاوتی داشته باشند.
	\item فرمت باینری تولیدی (مثلاً \lr{loader} و \lr{ABI}) ممکن است با آنچه کرنل روی برد انتظار دارد سازگار نباشد.
\end{itemize}
در نتیجه اجرای مستقیم چنین باینری‌ای روی برد یا اصلاً اجرا نمی‌شود (خطای \lr{Exec format error}) یا در بهترین حالت هنگام \lr{link} شدن پویا با کتابخانه‌های نامناسب دچار خطا یا رفتار نادرست می‌شود.

تفاوت اصلی: در \lr{Native Compilation}، معماری و محیط اجرای \lr{Host} (جایی که کامپایل انجام می‌شود) با معماری و محیط اجرای \lr{Target} (جایی که برنامه اجرا می‌شود) یکسان است. در \lr{Cross Compilation}، این دو محیط متفاوت‌اند: کامپایلر روی \lr{Host} اجرا می‌شود اما باینری‌ای تولید می‌کند که فقط روی \lr{Target} (معماری، \lr{ABI} و \lr{C Library} متفاوت) قابل اجراست. \lr{Host} معمولاً پردازنده‌ی قدرتمندتر توسعه‌دهنده (مثلاً \lr{x86}) و \lr{Target} پردازنده‌ی نهایی سیستم نهفته (مثلاً \lr{ARM}) است.

\subsubsection*{ب) اجزای \lr{Toolchain} و اهمیت ابزارهای \lr{Build}}

یک \lr{Cross Toolchain} برای \lr{C++} معمولاً شامل اجزای زیر است:
\begin{itemize}
	\item \textbf{\lr{Cross Compiler} (مثلاً \lr{arm-linux-gnueabihf-g++}):} کد \lr{C++} را به کد اسمبلی/شیء (\lr{object code}) مخصوص معماری \lr{ARM} ترجمه می‌کند.
	\item \textbf{\lr{Binutils} (\lr{assembler}، \lr{linker}، \lr{archiver}):} \lr{assembler} کد اسمبلی را به کد ماشین تبدیل می‌کند و \lr{linker} فایل‌های شیء و کتابخانه‌ها را با هم ترکیب کرده و باینری نهایی قابل اجرا روی \lr{Target} را می‌سازد.
	\item \textbf{\lr{Sysroot}:} مجموعه‌ای از هدرها و کتابخانه‌های مخصوص \lr{Target} (شامل \lr{C Library} و کتابخانه‌های \lr{C++} استاندارد نسخه‌ی \lr{Target}) که کامپایلر و \lr{linker} برای یافتن و \lr{link} کردن صحیح در برابر آن‌ها استفاده می‌کنند، نه نسخه‌های نصب‌شده روی \lr{Host}.
	\item \textbf{\lr{C/C++ Runtime Library}} (مانند \lr{glibc} یا \lr{musl} برای \lr{Target}): پیاده‌سازی توابع استاندارد و رابط \lr{system call} که برنامه در زمان اجرا روی \lr{Target} به آن نیاز دارد.
	\item \textbf{ابزارهای کمکی} مانند \lr{objdump}، \lr{readelf}، \lr{gdb} مخصوص \lr{Target} برای بازرسی و اشکال‌زدایی باینری تولیدشده.
\end{itemize}

با بزرگ‌تر شدن پروژه، تعداد فایل‌های منبع، وابستگی‌های بین آن‌ها، و گزینه‌های کامپایل (فلگ‌های \lr{Cross Compiler}، مسیرهای \lr{include}، کتابخانه‌های خارجی) به‌سرعت افزایش می‌یابد. فراخوانی دستی \lr{g++} برای هر فایل، هم پرخطا و هم غیرقابل نگهداری می‌شود. ابزارهایی مانند \lr{Make} با تعریف وابستگی بین فایل‌ها و اجرای فقط مراحل لازم (\lr{incremental build})، و \lr{CMake} با تولید خودکار فایل‌های \lr{build} برای \lr{toolchain}های مختلف (از جمله پیکربندی \lr{cross compilation} از طریق \lr{toolchain file})، مدیریت این پیچیدگی را ممکن می‌سازند و امکان \lr{build} یکسان و تکرارپذیر پروژه روی ماشین‌های مختلف را فراهم می‌کنند.

%% ────────────────────────────────────────────────────────────
\subsection{سؤال ۳: \lr{Root File System} و فایل‌سیستم‌ها در \lr{Linux Embedded}}

\subsubsection*{الف) پوشه‌های اصلی \lr{Root File System}}

\begin{itemize}
	\item \textbf{\lr{/bin}:} باینری‌های اجرایی ضروری سیستم (مانند \lr{ls}، \lr{sh}) که برای بوت و کار عادی سیستم حتی در حالت تک‌کاربره لازم‌اند.
	\item \textbf{\lr{/etc}:} فایل‌های پیکربندی سیستم و سرویس‌ها (مانند فایل‌های \lr{unit} سرویس \lr{systemd}، تنظیمات شبکه).
	\item \textbf{\lr{/lib}:} کتابخانه‌های اشتراکی (\lr{shared libraries}) مورد نیاز باینری‌های \lr{/bin} و \lr{/sbin}، از جمله \lr{C Library} و \lr{loader} پویا.
	\item \textbf{\lr{/dev}:} فایل‌های دستگاه (\lr{device nodes}) که رابط دسترسی به سخت‌افزار از طریق فایل‌سیستم را فراهم می‌کنند (معمولاً به‌صورت پویا توسط \lr{devtmpfs}/\lr{udev} ساخته می‌شوند).
	\item \textbf{\lr{/proc}:} فایل‌سیستم مجازی حاوی اطلاعات لحظه‌ای درباره‌ی پردازش‌ها و وضعیت کرنل.
	\item \textbf{\lr{/sys}:} فایل‌سیستم مجازی که مدل دستگاه‌ها و درایورهای کرنل (\lr{device model}) را به‌صورت سلسله‌مراتبی در معرض دید کاربر قرار می‌دهد.
	\item \textbf{\lr{/var}:} داده‌های متغیر در زمان اجرا مانند لاگ‌ها، فایل‌های موقت و صف‌ها.
\end{itemize}

تفاوت \lr{/proc} و \lr{/sys} با پوشه‌های معمولی: این دو، فایل‌سیستم‌های مجازی (\lr{virtual/pseudo filesystems}) هستند و هیچ محتوایی روی حافظه‌ی دائمی (فلش/دیسک) ندارند. فایل‌ها و پوشه‌های داخل آن‌ها در لحظه‌ی خواندن، توسط کرنل و از روی ساختارهای داده‌ی داخلی آن (مانند جدول پردازش‌ها یا مدل دستگاه‌ها) به‌صورت پویا تولید می‌شوند؛ نوشتن در برخی از این فایل‌ها نیز معادل تغییر مستقیم پارامترهای کرنل در حافظه است، نه ذخیره‌سازی یک فایل واقعی. به همین دلیل حجم آن‌ها صفر گزارش می‌شود و محتوایشان با ری‌استارت سیستم از نو ساخته می‌شود.

\subsubsection*{ب) مقایسه‌ی \lr{ext4}،‏ \lr{SquashFS} و \lr{tmpfs}}

\begin{table}[H]
\centering
\caption{مقایسه‌ی فایل‌سیستم‌های رایج در \lr{Linux Embedded}}
\begin{tabular}{lccc}
\toprule
\textbf{ویژگی} & \textbf{\lr{ext4}} & \textbf{\lr{SquashFS}} & \textbf{\lr{tmpfs}} \\
\midrule
نوع دسترسی        & \lr{Read-Write} & \lr{Read-Only} & \lr{Read-Write} \\
محل ذخیره‌سازی     & فلش/دیسک        & فلش/دیسک (فشرده) & \lr{RAM} \\
رفتار پس از قطع برق & نیاز به \lr{journaling} & بدون تغییر (فقط‌خواندنی) & از دست رفتن کامل داده \\
اثر بر \lr{wear} فلش & نوشتن مکرر، \lr{wear} بیشتر & فقط خوانده می‌شود، \lr{wear} صفر & بدون اثر (روی \lr{RAM}) \\
کاربرد مناسب        & پارتیشن داده‌ی پایدار (\lr{/data}) & \lr{Root FS} فشرده‌ی فقط‌خواندنی & \lr{/tmp}, \lr{/var/run} \\
\bottomrule
\end{tabular}
\end{table}

\begin{itemize}
	\item \textbf{\lr{ext4}:} فایل‌سیستم ژورنالی (\lr{journaling}) با قابلیت خواندن/نوشتن کامل. \lr{journaling} باعث می‌شود پس از قطع ناگهانی برق، فایل‌سیستم بتواند وضعیت سازگار خود را بازیابی کند، اما نوشتن مکرر متادیتا و ژورنال باعث افزایش \lr{wear} روی حافظه‌ی فلش می‌شود. مناسب برای بخش‌هایی از سیستم که داده‌ی کاربر یا پیکربندی باید به‌صورت پایدار و قابل‌تغییر ذخیره شود (مثل یک پارتیشن \lr{/data}).
	\item \textbf{\lr{SquashFS}:} فایل‌سیستم فشرده و فقط‌خواندنی. چون هیچ نوشتنی روی آن انجام نمی‌شود، در برابر قطع برق کاملاً ایمن است (چیزی برای خراب شدن وجود ندارد) و هیچ \lr{wear}ای روی فلش تحمیل نمی‌کند؛ همچنین به دلیل فشرده‌سازی، فضای کمتری اشغال می‌کند. بسیار مناسب برای \lr{Root File System} اصلی سیستم‌های نهفته‌ای که نیازی به تغییر باینری‌ها/برنامه‌ها در زمان اجرا ندارند.
	\item \textbf{\lr{tmpfs}:} فایل‌سیستمی که به‌طور کامل در \lr{RAM} قرار دارد. سرعت خواندن/نوشتن بسیار بالاست و هیچ اثری روی فلش ندارد (چون اصلاً روی فلش ذخیره نمی‌شود)، اما با خاموش شدن یا ری‌استارت سیستم تمام محتوای آن از بین می‌رود. مناسب برای دایرکتوری‌هایی مانند \lr{/tmp} یا \lr{/var/run} که محتوای آن‌ها موقتی است و نیازی به ماندگاری ندارد؛ استفاده از آن به‌جای \lr{ext4} برای این مسیرها، \lr{wear} فلش را کاهش می‌دهد.
\end{itemize}

یک الگوی رایج در \lr{Linux Embedded}، ترکیب این سه است: \lr{Root File System} اصلی روی \lr{SquashFS} (فقط‌خواندنی، مقاوم در برابر قطع برق)، یک پارتیشن جداگانه‌ی \lr{ext4} برای داده‌های پایدار قابل‌تغییر، و \lr{tmpfs} برای مسیرهای موقتی پرنوشت.

%% ────────────────────────────────────────────────────────────
\subsection{سؤال ۴: \lr{Device Tree} و \lr{Device Tree Overlay}}

\subsubsection*{الف) چرا سخت‌افزار با \lr{Device Tree} توصیف می‌شود؟}

در گذشته، اطلاعات سخت‌افزاری بردهای \lr{ARM} به‌صورت \lr{hard-code} در کد کرنل (\lr{board files}) نوشته می‌شد؛ به این معنا که برای هر برد جدید یا حتی تغییر جزئی سخت‌افزار (مثلاً جابه‌جایی یک پین)، باید کد کرنل تغییر کرده و کرنل دوباره کامپایل می‌شد. \lr{Device Tree} این مشکل را با جدا کردن «توصیف سخت‌افزار» از «کد کرنل» حل می‌کند: یک فایل داده‌ای مستقل (\lr{.dts}) که در زمان بوت توسط \lr{Bootloader} به کرنل تحویل داده می‌شود، ساختار سخت‌افزار برد را توصیف می‌کند و کرنل بر اساس آن، درایورهای مناسب را با پارامترهای درست فعال می‌کند؛ بدون آنکه نیازی به تغییر یا کامپایل مجدد کد کرنل باشد.

\lr{Device Tree} معمولاً اطلاعاتی مانند موارد زیر را توصیف می‌کند:
\begin{itemize}
	\item نقشه‌ی حافظه (آدرس و اندازه‌ی \lr{RAM}، نواحی \lr{reserved})
	\item آدرس‌های ثبات‌ها (\lr{memory-mapped I/O}) و شماره‌ی وقفه‌ (\lr{IRQ}) هر پیرامونی
	\item رشته‌های \lr{compatible} که مشخص می‌کنند کدام درایور کرنل باید برای یک دستگاه بارگذاری شود
	\item پیکربندی گذرگاه‌های \lr{I2C}،‏ \lr{SPI}،‏ \lr{UART} و دستگاه‌های متصل به آن‌ها
	\item تنظیمات \lr{pinmux/GPIO} و درخت کلاک‌ها (\lr{clock tree})
\end{itemize}

\subsubsection*{ب) \lr{Device Tree Overlay}}

\lr{Device Tree Overlay} یک قطعه‌ی جداگانه از توصیف \lr{Device Tree} است که در زمان بوت (یا حتی در زمان اجرا) بر روی \lr{Device Tree} پایه‌ی برد \lr{اعمال (merge)} می‌شود، بدون آنکه نیاز به تغییر یا کامپایل مجدد \lr{Device Tree} اصلی یا کد کرنل باشد. این ویژگی زمانی مفید است که سخت‌افزار جانبی به‌صورت ماژولار و قابل‌تعویض به برد اصلی اضافه می‌شود -- مانند بردهای الحاقی \lr{HAT} روی \lr{Raspberry Pi} یا \lr{cape} روی \lr{BeagleBone} -- و توسعه‌دهنده می‌خواهد بسته به این‌که چه سخت‌افزاری فیزیکاً نصب شده، پیکربندی متفاوتی اعمال کند.

\textbf{مثال مشخص:} فرض کنید یک سنسور دما روی گذرگاه \lr{I2C1} به یک \lr{Raspberry Pi} متصل می‌شود. به‌جای ویرایش \lr{Device Tree} اصلی برد و کامپایل مجدد کرنل، می‌توان یک \lr{Overlay} کوچک نوشت که گره‌ای جدید زیر گره‌ی \lr{i2c1} در \lr{Device Tree} پایه اضافه می‌کند و مشخص می‌سازد که یک دستگاه با \lr{compatible = "bosch,bme280"} روی آدرس \lr{0x76} این گذرگاه قرار دارد. این \lr{Overlay} به‌صورت یک فایل \lr{.dtbo} کامپایل شده و فقط با افزودن نام آن به فایل پیکربندی \lr{Bootloader} (مثلاً \lr{config.txt} در \lr{Raspberry Pi}) در بوت بعدی فعال می‌شود؛ کرنل با دیدن این گره‌ی جدید، درایور مربوط به سنسور \lr{BME280} را بدون هیچ تغییری در کد یا \lr{Device Tree} اصلی بارگذاری می‌کند.

%% ────────────────────────────────────────────────────────────
\subsection{سؤال ۵: پیش‌بینی‌پذیری زمانی در \lr{Linux} و سیستم‌های بلادرنگ}

\subsubsection*{الف) چرا \lr{Linux} استاندارد برای \lr{Hard Real-Time} مناسب نیست؟}

هسته‌ی \lr{Linux} استاندارد برای \lr{throughput} و انصاف (\lr{fairness}) میان پردازش‌ها طراحی شده، نه برای تضمین قطعی زمان پاسخ. منابع اصلی \lr{latency} و \lr{jitter} در آن عبارت‌اند از:
\begin{itemize}
	\item \textbf{بخش‌های غیرقابل‌قطع کرنل (\lr{non-preemptible sections}):} بسیاری از مسیرهای کد کرنل (مانند \lr{spinlock}ها یا \lr{interrupt handler}های کوتاه) اجازه نمی‌دهند یک \lr{task} با اولویت بالاتر آن‌ها را قطع کند، حتی اگر آن \lr{task} یک \lr{deadline} سخت داشته باشد.
	\item \textbf{وقفه‌ها (\lr{interrupts}) و \lr{Interrupt Handler}ها:} پردازش وقفه‌ها می‌تواند اجرای یک \lr{task} بلادرنگ را برای مدتی نامشخص به تأخیر بیندازد، به‌خصوص وقتی چند دستگاه هم‌زمان وقفه تولید می‌کنند.
	\item \textbf{زمان‌بند (\lr{Scheduler}) و اولویت‌های نرم:} زمان‌بند پیش‌فرض \lr{Linux} (\lr{CFS}) برای انصاف طراحی شده، نه برای تضمین سخت \lr{deadline}؛ حتی با اولویت‌های \lr{real-time} (\lr{SCHED\_FIFO}/\lr{SCHED\_RR})، تأخیرهای ناشی از بخش‌های غیرقابل‌قطع همچنان وجود دارند.
	\item \textbf{مدیریت حافظه و \lr{cache}:} \lr{paging}، \lr{page fault}ها، و رفتار غیرقطعی \lr{cache} (بسته به الگوی دسترسی سایر پردازش‌ها) باعث تغییرات غیرقابل‌پیش‌بینی در زمان اجرای یک قطعه کد می‌شوند.
	\item \textbf{قفل‌های غیرقطعی (\lr{non-deterministic locking}):} انتظار برای گرفتن \lr{mutex}هایی که ممکن است توسط پردازش‌های با اولویت پایین‌تر نگه داشته شده باشند (\lr{priority inversion}) نیز منبع دیگری از \lr{jitter} است.
\end{itemize}

\subsubsection*{ب) \lr{PREEMPT\_RT} در برابر واگذاری به \lr{MCU}/\lr{RTOS} جداگانه}

\begin{table}[H]
\centering
\caption{مقایسه‌ی دو رویکرد دستیابی به رفتار زمانی قابل پیش‌بینی}
\begin{tabular}{p{4cm}p{5.3cm}p{5.3cm}}
\toprule
 & \textbf{\lr{Linux + PREEMPT\_RT}} & \textbf{\lr{MCU/RTOS} جداگانه} \\
\midrule
مزیت اصلی &
یک محیط برنامه‌نویسی واحد (\lr{Linux API}ها، \lr{driver}ها، \lr{networking}) با \lr{latency} بسیار کاهش‌یافته &
جداسازی کامل بخش زمان‌حساس از پیچیدگی \lr{Linux}؛ \lr{worst-case latency} در حد میکروثانیه و کاملاً قطعی \\
محدودیت &
همچنان \lr{soft/firm real-time} است، نه \lr{hard} به‌معنای واقعی؛ \lr{worst-case latency} در حد ده‌ها میکروثانیه تا میلی‌ثانیه باقی می‌ماند و به سخت‌افزار/بار سیستم وابسته است &
نیاز به دو محیط توسعه، دو \lr{toolchain}، و یک لایه‌ی ارتباطی (مثلاً \lr{shared memory},‏ \lr{UART}, یا \lr{SPI}) بین \lr{MCU} و \lr{Linux} برای هماهنگی \\
\lr{Trade-off} &
پیچیدگی توسعه کم، اما تضمین زمانی محدود &
تضمین زمانی قوی، اما پیچیدگی معماری و توسعه‌ی سیستم بیشتر \\
مناسب برای &
سیستم‌هایی با \lr{deadline}های نسبتاً نرم تا سخت با تلورانس چند ده میکروثانیه (مثلاً کنترل صنعتی با نرخ متوسط، \lr{data acquisition}) &
سیستم‌هایی با \lr{deadline}های بسیار سخت و بحرانی (مثلاً کنترل موتور با حلقه‌ی بسته‌ی بسیار سریع، ایمنی خودرو) که هیچ‌گونه تأخیر غیرمنتظره قابل‌قبول نیست \\
\bottomrule
\end{tabular}
\end{table}

به‌طور کلی، اگر نیاز سیستم، اجرای منطق پیچیده (پردازش تصویر، شبکه، رابط کاربری) در کنار محدودیت‌های زمانی نسبتاً قابل‌تحمل باشد، \lr{PREEMPT\_RT} انتخاب ساده‌تر و کم‌هزینه‌تری است. اما اگر بخشی از سیستم واقعاً به \lr{deadline}های سخت و قطعی (میکروثانیه‌ای) نیاز داشته باشد -- جایی که حتی یک تأخیر نادر می‌تواند فاجعه‌بار باشد -- واگذاری آن بخش به یک \lr{microcontroller}/\lr{RTOS} مجزا و استفاده از \lr{Linux} تنها برای بخش‌های غیر‌بحرانی (رابط کاربری، \lr{logging}، ارتباطات شبکه) انتخاب مطمئن‌تری است.


%% ============================================================
\section{بخش عملی}
%% ============================================================

\subsection{تمرین عملی ۱: سرویس محلی درخواست و پاسخ با \lr{Unix Domain Socket}}

\subsubsection*{طراحی پروتکل}

یک \lr{Unix Domain Socket} از نوع \lr{SOCK\_STREAM} روی مسیر \lr{/tmp/hw3\_ex1.sock} استفاده شده است. پروتکل، یک پروتکل متنی ساده و خط‌محور (\lr{newline-terminated}) است:

\begin{itemize}
	\item هر درخواست، یک خط متنی است که با \lr{`\textbackslash n`} پایان می‌یابد.
	\item فرمان \lr{DATE}: بدون آرگومان؛ سرور تاریخ و زمان فعلی سیستم را برمی‌گرداند.
	\item فرمان \lr{ECHO <text>}: باقی‌مانده‌ی خط پس از کلمه‌ی \lr{ECHO} (که می‌تواند خود شامل فاصله باشد) عیناً به‌عنوان پاسخ بازگردانده می‌شود.
	\item پاسخ سرور نیز یک خط است: در صورت موفقیت با پیشوند \lr{OK}، و در صورت نامعتبر بودن درخواست با پیشوند \lr{ERR}.
\end{itemize}

با توجه به این‌که در صورت مسئله سرویس‌دهی باید «پشت سر هم» به چند \lr{client} انجام شود (نه لزوماً هم‌زمان)، طراحی به‌صورت یک حلقه‌ی \lr{accept} تکراری (\lr{iterative}, بدون \lr{thread}/\lr{fork}) پیاده‌سازی شده است: هر اتصال یک درخواست/پاسخ را رد و بدل می‌کند، سپس بسته می‌شود و سرور به پذیرش اتصال بعدی بازمی‌گردد. این طراحی از پیچیدگی همگام‌سازی (\lr{synchronization}) میان \lr{thread}ها بی‌نیاز است و کاملاً با عبارت «سرویس‌دهی پشت سر هم» در صورت مسئله مطابقت دارد؛ در صورت نیاز به پاسخ‌گویی هم‌زمان به چند \lr{client}، گام بعدی طبیعی استفاده از یک \lr{thread} به ازای هر اتصال یا \lr{multiplexing} با \lr{poll}/\lr{epoll} خواهد بود.

خطاهای راه‌اندازی سوکت (\lr{socket}/\lr{bind}/\lr{listen}) باعث خروج برنامه با پیام خطا می‌شوند، اما خطای مربوط به یک اتصال خاص (مثلاً قطع ناگهانی \lr{client}) فقط باعث بسته شدن همان اتصال می‌شود و حلقه‌ی اصلی سرور را متوقف نمی‌کند.

\subsubsection*{کد سرور}

\begin{latin}
\lstinputlisting[language=C++, caption={\rl{server.cpp -- سرور \lr{Unix Domain Socket}}}]{../code/ex1_socket/server.cpp}
\end{latin}

\subsubsection*{کد کلاینت}

\begin{latin}
\lstinputlisting[language=C++, caption={\rl{client.cpp -- کلاینت \lr{Unix Domain Socket}}}]{../code/ex1_socket/client.cpp}
\end{latin}

\subsubsection*{\lr{Makefile}}

\begin{latin}
\lstinputlisting[language=make, caption={\rl{Makefile -- ساخت سرور و کلاینت}}]{../code/ex1_socket/Makefile}
\end{latin}

\subsubsection*{نتایج اجرا}

نحوه‌ی اجرا: \lr{./server} در یک ترمینال اجرا و روی سوکت گوش می‌دهد؛ در ترمینال دوم دستور \lr{./client date} یا \lr{./client echo <text>} اجرا می‌شود.

همان‌طور که در شکل زیر دیده می‌شود، هر دو فرمان \lr{DATE} و \lr{ECHO} در یک نشست از \lr{client} روی سرور اوبونتو آزموده شده‌اند؛ لاگ سمت سرور (ترمینال چپ) نیز درخواست و پاسخ متناظر با هرکدام را نشان می‌دهد.

\begin{figure}[H]
\centering
\includegraphics[width=\textwidth]{ex1_date_demo.png}
\caption{اجرای هم‌زمان سرور (چپ) و کلاینت (راست): تبادل فرمان‌های DATE و ECHO}
\label{fig:ex1-demo}
\end{figure}

%% ────────────────────────────────────────────────────────────
\subsection{تمرین عملی ۲: پایشگر محلی سیستم (\lr{procwatch})}

برنامه‌ی \lr{procwatch} با دو حالت اجرای \lr{interactive} و \lr{background} پیاده‌سازی شده و کد آن در پوشه‌های \lr{src} و \lr{include} سازمان‌دهی شده است (به \lr{README.md} همراه کد برای دستور دقیق ساخت و اجرا مراجعه شود).

\subsubsection*{بخش ۱: اطلاعات کلی سیستم}

نسخه‌ی کرنل از طریق فراخوانی \lr{uname()} (به‌جای پارس مستقیم \lr{/proc/version}، هرچند هر دو معتبرند) و تعداد هسته‌ها با شمارش خطوط \lr{processor} در \lr{/proc/cpuinfo} به دست می‌آید (معادل \lr{sysconf(\_SC\allowbreak\_NPROCESSORS\allowbreak\_ONLN)}).

\textbf{درصد استفاده‌ی کلی \lr{CPU}:} خط اول فایل \lr{/proc/stat} مقادیر تجمعی (بر حسب \lr{jiffies}/\lr{USER\_HZ}) شمارنده‌ی زمان صرف‌شده در هر حالت را از لحظه‌ی بوت سیستم نشان می‌دهد:
\[
\text{\lr{cpu}} \quad \text{\lr{user}}\ \ \text{\lr{nice}}\ \ \text{\lr{system}}\ \ \text{\lr{idle}}\ \ \text{\lr{iowait}}\ \ \text{\lr{irq}}\ \ \text{\lr{softirq}}\ \ \text{\lr{steal}}
\]
در این تمرین از مجموع تمام این فیلدها به‌عنوان \lr{total} و از مجموع \lr{idle + iowait} به‌عنوان \lr{idle} استفاده شده است. از آنجا که این مقادیر شمارنده‌های \emph{تجمعی از لحظه‌ی بوت} هستند (نه نرخ لحظه‌ای)، یک بار خواندن این فایل هیچ اطلاعاتی درباره‌ی میزان اشغال \lr{CPU} در «حال حاضر» به ما نمی‌دهد -- فقط می‌گوید از ابتدای بوت چه مقدار زمان صرف چه کاری شده. برای به‌دست‌آوردن درصد استفاده در یک بازه‌ی زمانی مشخص، باید دو نمونه با فاصله‌ی زمانی معلوم (در این پیاده‌سازی ۲ ثانیه، هم‌زمان با دوره‌ی به‌روزرسانی جدول) گرفته و از تفاضل آن‌ها نرخ اشغال را محاسبه کرد:
\[
\%\text{\lr{busy}} = 100 \times \left(1 - \frac{\Delta \text{\lr{idle}}}{\Delta \text{\lr{total}}}\right)
\]
این فرمول در تابع \lr{cpu\_percent} در فایل \lr{system\_info.cpp} پیاده‌سازی شده است.

حافظه از \lr{/proc/meminfo} (فیلدهای \lr{MemTotal}،‏ \lr{MemFree}،‏ \lr{MemAvailable}، همگی بر حسب \lr{kB}) و \lr{load average} مستقیماً از سه عدد اول \lr{/proc/loadavg} خوانده می‌شوند (این مقدار توسط خود کرنل به‌صورت میانگین متحرک نمایی محاسبه شده و نیازی به نمونه‌گیری مضاعف ندارد).

\begin{latin}
\lstinputlisting[language=C++, caption={\rl{system\_info.cpp -- اطلاعات کلی سیستم و محاسبه‌ی \lr{CPU\%}}}]{../code/ex2_procwatch/src/system_info.cpp}
\end{latin}

\subsubsection*{بخش ۲: پردازش‌های پرمصرف (\lr{Top-N})}

پردازش‌ها با پیمایش پوشه‌های عددیِ \lr{/proc} (از طریق \lr{opendir}/\lr{readdir}، فیلتر بر اساس این‌که نام پوشه کاملاً عددی باشد) شناسایی می‌شوند. برای هر \lr{PID}، فایل \lr{/proc/[pid]/stat} خوانده می‌شود. از آن‌جا که نام پردازش (\lr{comm}) می‌تواند خود شامل فاصله یا پرانتز باشد، پارس این فایل با یافتن \textbf{آخرین} کاراکتر \lr{`)'} در خط (نه اولین) انجام می‌شود؛ محتوای بین اولین \lr{`('} و این آخرین \lr{`)'} نام پردازش است و بقیه‌ی فیلدها با تفکیک بر اساس فاصله، به‌صورت موقعیتی خوانده می‌شوند: فیلد چهاردهم \lr{utime} و پانزدهم \lr{stime} (هر دو بر حسب \lr{clock ticks}).

میزان فعالیت \lr{CPU} هر پردازش از تفاضل \lr{utime + stime} بین دو نمونه‌ی متوالی (با همان فاصله‌ی ۲ ثانیه‌ی به‌روزرسانی) به دست می‌آید. طبق صورت مسئله نیازی به تبدیل این مقدار به درصد واقعی نیست، اما برای این تبدیل باید بر حاصل‌ضرب \lr{CLK\_TCK} (تعداد \lr{ticks} در ثانیه، قابل‌دریافت از \lr{sysconf(\_SC\_CLK\_TCK)}، معمولاً برابر ۱۰۰ در لینوکس) در طول بازه‌ی زمانی نمونه‌گیری برحسب ثانیه تقسیم شود:
\[
\%\text{\lr{cpu}}_{\text{process}} = 100 \times \frac{\Delta(\text{\lr{utime}}+\text{\lr{stime}})}{\text{\lr{CLK\_TCK}} \times \text{\lr{interval\_seconds}}}
\]

مقدار \lr{RSS} از فایل \lr{/proc/[pid]/status} و خط \lr{VmRSS} خوانده می‌شود که از قبل بر حسب \lr{kB} گزارش می‌شود (برخلاف فیلد \lr{rss} در \lr{/proc/[pid]/stat} که بر حسب تعداد صفحه‌ی حافظه است و نیاز به ضرب در \lr{sysconf(\_SC\_PAGESIZE)} دارد)؛ استفاده از \lr{VmRSS} ساده‌تر است و در این پیاده‌سازی انتخاب شده است.

برای مقاوم بودن در برابر خطاهای رایج \lr{/proc} (مانند پایان یافتن یک پردازش میان زمان \lr{readdir} و زمان باز کردن فایل آن، که منجر به خطاهایی نظیر \lr{ENOENT}/\lr{ESRCH} می‌شود)، هرگاه باز کردن یا خواندن فایل \lr{/proc/[pid]/stat} یا \lr{/proc/[pid]/status} ناموفق باشد، آن \lr{PID} به‌سادگی از این دوره‌ی نمونه‌گیری صرف‌نظر می‌شود، بدون آنکه برنامه متوقف یا دچار \lr{crash} شود.

در نهایت، پردازش‌های موجود در هر دو نمونه (قبلی و فعلی) بر اساس میزان فعالیت \lr{CPU} به‌صورت نزولی مرتب شده و \lr{N} پردازش برتر (پیش‌فرض \lr{N=5}) انتخاب می‌شوند.

\begin{latin}
\lstinputlisting[language=C++, caption={\rl{proc\_scanner.cpp -- پیمایش \lr{/proc} و محاسبه‌ی \lr{Top-N}}}]{../code/ex2_procwatch/src/proc_scanner.cpp}
\end{latin}

\begin{latin}
\lstinputlisting[language=C++, caption={\rl{interactive\_mode.cpp -- نمایش زنده‌ی جدول‌ها}}]{../code/ex2_procwatch/src/interactive_mode.cpp}
\end{latin}

\subsubsection*{بخش ۳: اجرای برنامه به‌صورت سرویس \lr{systemd}}

در حالت \lr{background}، برنامه با \lr{procwatch background <config-file>} اجرا می‌شود. فایل پیکربندی یک فرمت ساده‌ی \lr{key=value} است که حداقل \lr{interval} (فاصله‌ی زمانی بر حسب ثانیه بین ثبت‌های \lr{Top-N} در \lr{journal}) و \lr{N} (تعداد پردازش‌های برتر) را شامل می‌شود؛ همچنین یک فیلد \lr{fail} برای آزمون مسیر خطای عمدی تعبیه شده است.

از آن‌جا که \lr{systemd} به‌صورت پیش‌فرض خروجی استاندارد (\lr{stdout}/\lr{stderr}) یک سرویس را در \lr{journal} ثبت می‌کند، در این پیاده‌سازی نیازی به فراخوانی مستقیم \lr{API} کتابخانه‌ی \lr{libsystemd} نیست؛ کافی است هر دوره ثبت‌ها با \lr{std::cout}/\lr{std::cerr} چاپ و بلافاصله \lr{flush} شوند (چون خروجی غیر\lr{TTY} به‌صورت \lr{buffered} است و بدون \lr{flush} ممکن است دیرتر در \lr{journal} ظاهر شود).

\textbf{تفاوت \lr{SIGTERM} و خطای عمدی:} یک \lr{signal handler} برای \lr{SIGTERM} نصب شده که فقط یک پرچم \lr{volatile sig\_atomic\_t} را تنظیم می‌کند (برای رعایت \lr{async-signal-safety}؛ هیچ‌گونه چاپ یا تخصیص حافظه داخل خود \lr{handler} انجام نمی‌شود). حلقه‌ی اصلی برنامه این پرچم را در هر تکرار بررسی می‌کند و در صورت فعال بودن، پیام پایان‌کار را ثبت کرده و با کد \lr{0} خارج می‌شود. در مقابل، اگر مقدار \lr{fail=true} در فایل پیکربندی تنظیم شده باشد، برنامه بلافاصله در ابتدای اجرا پیام خطا را ثبت کرده و با کد خروج \lr{1} (غیرصفر) متوقف می‌شود -- این مسیر برای آزمودن رفتار \lr{Restart=on-failure} سرویس \lr{systemd} استفاده می‌شود.

\begin{latin}
\lstinputlisting[language=C++, caption={\rl{background\_mode.cpp -- حالت سرویس، پیکربندی و مدیریت \lr{SIGTERM}}}]{../code/ex2_procwatch/src/background_mode.cpp}
\end{latin}

\textbf{فایل \lr{unit}} (\lr{procwatch.service}):

\begin{latin}
\lstinputlisting[caption={\rl{procwatch.service -- تعریف سرویس \lr{systemd}}}]{../code/ex2_procwatch/procwatch.service}
\end{latin}

فیلد \lr{Type=simple} به این معناست که فرآیند اصلی \lr{ExecStart} همان فرآیند سرویس است (بدون \lr{fork} کردن دیمن)؛ \lr{Restart=on-failure} به همراه \lr{RestartSec=5} باعث می‌شود سرویس \textbf{فقط} در صورت خروج با کد غیرصفر یا کشته‌شدن با یک سیگنال غیرمنتظره، پس از حدود ۵ ثانیه دوباره اجرا شود -- در حالی که یک خروج تمیز پس از \lr{SIGTERM} (کد ۰) هیچ راه‌اندازی مجددی را فعال نمی‌کند. \lr{WantedBy=multi-user.target} به همراه \lr{systemctl enable} سرویس را در زمان بوت سیستم فعال می‌کند.

\textbf{فایل پیکربندی نمونه} (\lr{procwatch.conf}):

\begin{latin}
\lstinputlisting[caption={\rl{procwatch.conf -- پیکربندی نمونه‌ی حالت \lr{background}}}]{../code/ex2_procwatch/procwatch.conf}
\end{latin}

\subsubsection*{نتایج اجرا}

\begin{figure}[H]
\centering
\includegraphics[width=0.85\textwidth]{ex2_make.png}
\caption{ساخت procwatch با Make}
\label{fig:ex2-build-make}
\end{figure}

\begin{figure}[H]
\centering
\includegraphics[width=0.85\textwidth]{ex2_cmake.png}
\caption{ساخت procwatch با CMake}
\label{fig:ex2-build-cmake}
\end{figure}

\textbf{یادداشت اصلاحیه:} در بازبینی اسکرین‌شات‌های اولیه‌ی اجرا، مشخص شد ستون \lr{RSS(KB)} همواره مقدار صفر نشان می‌دهد. علت، اشکالی در تابع خواندن \lr{VmRSS} از \lr{/proc/[pid]/status} بود: پیاده‌سازی اولیه این فایل را با الگوی تکرارشونده‌ی سه‌تایی \lr{key >> value >> unit} می‌خواند، اما فیلدهایی مانند \lr{Umask:} یا \lr{Uid:} که پیش از \lr{VmRSS:} در این فایل قرار دارند، دقیقاً این قالب سه‌تایی را ندارند و باعث می‌شوند این روش خواندن، هم‌ترازی خود با ابتدای هر سطر را از دست بدهد و هرگز به‌درستی به سطر \lr{VmRSS:} نرسد. این اشکال با تغییر پیاده‌سازی به خواندن خط‌به‌خط و تطبیق مستقیم پیشوند \lr{"VmRSS:"} در تابع \lr{parse\_rss\_kb} (فایل \lr{proc\_scanner.cpp}) برطرف شد؛ کد بالا در بخش‌های قبلی همین گزارش، نسخه‌ی اصلاح‌شده را نشان می‌دهد. اسکرین‌شات‌های زیر پس از بازساخت برنامه با این اصلاح، دوباره گرفته شده‌اند و اکنون مقادیر \lr{RSS} درست (غیرصفر برای پردازش‌های دارای حافظه‌ی مقیم) را نشان می‌دهند؛ مقدار صفر برای برخی \lr{kworker}ها طبیعی است، چون این نخ‌های کرنلی صفحه‌ی حافظه‌ی کاربری مقیمی ندارند.

\begin{figure}[H]
\centering
\includegraphics[width=0.7\textwidth]{ex2_procwatch.png}
\caption{اجرای حالت interactive در حالت آرام سیستم، با مقادیر RSS صحیح}
\label{fig:ex2-interactive-idle}
\end{figure}

\begin{figure}[H]
\centering
\includegraphics[width=\textwidth]{ex2_stresstest.png}
\caption{ایجاد بار روی یک هسته با stress-ng هم‌زمان با اجرای procwatch interactive (چپ) و htop (راست) برای مقایسه}
\label{fig:ex2-load-test}
\end{figure}

\begin{figure}[H]
\centering
\includegraphics[width=0.85\textwidth]{ex2_status.png}
\caption{خروجی systemctl status procwatch پس از بازساخت با کد اصلاح‌شده، با مقادیر RSS صحیح در لاگ}
\label{fig:ex2-systemctl-status}
\end{figure}

\begin{figure}[H]
\centering
\includegraphics[width=0.85\textwidth]{ex2_journalctl.png}
\caption{لاگ‌های دوره‌ای procwatch در journalctl با مقادیر RSS صحیح}
\label{fig:ex2-journalctl}
\end{figure}

توقف صحیح سرویس با \lr{SIGTERM} پیش از اصلاح فوق آزموده شده و همچنان معتبر است (اشکال \lr{RSS} تأثیری بر رفتار خاموش‌شدن سرویس ندارد):

\begin{figure}[H]
\centering
\includegraphics[width=0.85\textwidth]{ex2_stop.png}
\caption{توقف صحیح سرویس procwatch با SIGTERM: پیام پایان‌کار در journal و وضعیت inactive (dead) با کد خروج 0}
\label{fig:ex2-sigterm-stop}
\end{figure}

به همین ترتیب، راه‌اندازی مجدد پس از خطای عمدی (\lr{fail=true}) نیز به رفتار \lr{RSS} وابسته نیست و نتیجه‌ی زیر همچنان معتبر است:

\begin{figure}[H]
\centering
\includegraphics[width=\textwidth]{ex2_restart.png}
\caption{چرخه‌ی failed $\to$ راه‌اندازی مجدد procwatch پس از خطای عمدی (fail=true)، با فاصله‌ی حدود ۵ ثانیه و افزایش شمارنده‌ی راه‌اندازی مجدد}
\label{fig:ex2-restart-failure}
\end{figure}
